% ID: SEL-D0065-OB-E02
% ESTADO: validada. S065.
\textbf{a)} El dominio es $\mathbb R\setminus\{0\}$. Fuera de la unión en $1$, las ramas son continuas y derivables en su dominio. En $1$ ambos límites coinciden con $f(1)=2$. Las pendientes laterales son $-2/x^2$ y $2x-4$ evaluadas en $1$, por lo que ambas valen $-2$. Así, la función es continua y derivable en todo su dominio. En $0$ tiene una asíntota vertical.

\textbf{b)} Para $x<0$, la hipérbola decrece desde valores próximos a cero hasta $-\infty$. Para $0<x\le1$ decrece desde $+\infty$ hasta $(1,2)$. La parábola $(x-2)^2+1$ sigue decreciendo hasta $(2,1)$ y luego crece. La unión en $(1,2)$ tiene tangente común.
\begin{center}
\setlength{\unitlength}{0.9mm}
\begin{picture}(120,70)(-53,-33)
\qbezier(-47,0)(7,0)(61,0)
\qbezier(0,-29)(0,1)(0,31)
\put(62,0){$x$}\put(2,31){$y$}
\qbezier(-45.6,-2.63158)(-40.8,-2.8999)(-36,-3.33333)
\qbezier(-36,-3.33333)(-30,-3.83333)(-24,-5)
\qbezier(-24,-5)(-19.5,-5.80769)(-15,-8)
\qbezier(-15,-8)(-12.3,-9.2622)(-9.6,-12.5)
\qbezier(-9.6,-12.5)(-8.1,-14.2887)(-6.6,-18.1818)
\qbezier(-6.6,-18.1818)(-5.7,-20.5144)(-4.8,-25)
\qbezier(4.8,25)(5.4,21.9444)(6,20)
\qbezier(6,20)(6.9,17.0903)(7.8,15.3846)
\qbezier(7.8,15.3846)(8.7,13.6439)(9.6,12.5)
\qbezier(9.6,12.5)(10.8,10.9722)(12,10)
\qbezier(12,10)(30,-5)(48,25)
\put(12,12){\scriptsize $(1,2)$}
\put(21,8){\scriptsize $(2,1)$}
\end{picture}
\end{center}
